\begin{problem}[33届物理竞赛决赛]{67}
    四、具有一定能量、动量的光子还具有角动量。圆偏振光的光子的角动量大小为 $\hbar$ 。光子被物体吸收后，光子的能量、动量和角动量就全部传给物体。物体吸收光子获得的角动量可以使物体转动。科学家利用这一原理，在连续的圆偏振激光照射下，实现了纳米颗粒的高速转动，获得了迄今为止液体环境中转速最高的微尺度转子。

    如图，一金纳米球颗粒放置在两片水平光滑玻璃平板之间，并整体（包括玻璃平板）浸在水中。一束圆偏振激光从上往下照射到金纳米颗粒上。已知该束入射激光在真空中的波长 $\lambda = 830\,\mathrm{nm}$ ，经显微物镜聚焦后（仍假设为平面波，每个光子具有沿传播方向的角动量 $\hbar$ ）光斑直径 $d = 1.20 \times 10^{-6}\,\mathrm{m}$ ，功率 $P = 20.0\,\mathrm{mW}$ 。金纳米球颗粒的半径 $R = 100\,\mathrm{nm}$ ，金的密度 $\rho = 19.32 \times 10^{3}\,\mathrm{kg/m^{3}}$ 。忽略光在介质界面上的反射以及玻璃、水对光的吸收等损失，仅从金纳米颗粒吸收光子获得角动量驱动其转动的角度分析下列问题（计算结果取3位有效数字）：
    \begin{subproblem}
        求该束激光的频率 $\nu$ 和光强 $I$ （在传播方向上单位横截面积所传输的功率）。
    \end{subproblem}
    \begin{subproblem}
        给出金纳米球颗粒的质量 $m$ 和它绕其对称轴的转动惯量 $J$ 的值。
    \end{subproblem}
    \begin{subproblem}
        假设颗粒对光的吸收截面（颗粒吸收的光功率与入射光强之比）为 $\sigma_{\mathrm{abs}}=0.123\pi R^{2}$ ，求该束激光作用在颗粒上沿旋转对称轴的力矩的大小 $M$。
    \end{subproblem}
    \begin{subproblem}
        已知一个以角速度为 $\omega$ 旋转的球形颗粒（半径为 $r$）在粘滞系数为 $\eta$ 的液体中受到的粘滞摩擦力矩大小为 $M_{f}=8\pi\eta r^{3}\omega$ 。已知水的粘滞系数 $\eta=8.00\times10^{-4}\,\mathrm{Pa\cdot s}$ ，求金纳米球颗粒在此光束照射下达到稳定旋转时的转速（转数/秒） $f$ 。
    \end{subproblem}
    \begin{subproblem}
        取光开始照到处于静止状态的金纳米颗粒的瞬间为计时零点 $t_{0}=0$ ，求在任意时刻 $t$ ($t>0$) 该颗粒转速的表达式 $f(t)$ 。
    \end{subproblem}
    \begin{subproblem}
        若把入射激光束换成方波脉冲激光束，脉冲宽度为 $T_{1}$ （此期间内光强仍为 $I$），脉冲之间的间歇时间为 $T_{2}$ 。取第一个脉冲的光开始照到颗粒的时刻为计时零点 $t_{0}=0$ ，求第 $n$ 个完整脉冲周期 $(T_{1}+T_{2})$ 后的瞬间颗粒的转速 $f_{n}$ 的表达式，并给出转速极限 $f_{n\to\infty}$ 的表达式。
    \end{subproblem}
\end{problem}